% Options for packages loaded elsewhere
\PassOptionsToPackage{unicode}{hyperref}
\PassOptionsToPackage{hyphens}{url}
\documentclass[
]{article}
\usepackage{xcolor}
\usepackage[margin=1in]{geometry}
\usepackage{amsmath,amssymb}
\setcounter{secnumdepth}{-\maxdimen} % remove section numbering
\usepackage{iftex}
\ifPDFTeX
  \usepackage[T1]{fontenc}
  \usepackage[utf8]{inputenc}
  \usepackage{textcomp} % provide euro and other symbols
\else % if luatex or xetex
  \usepackage{unicode-math} % this also loads fontspec
  \defaultfontfeatures{Scale=MatchLowercase}
  \defaultfontfeatures[\rmfamily]{Ligatures=TeX,Scale=1}
\fi
\usepackage{lmodern}
\ifPDFTeX\else
  % xetex/luatex font selection
\fi
% Use upquote if available, for straight quotes in verbatim environments
\IfFileExists{upquote.sty}{\usepackage{upquote}}{}
\IfFileExists{microtype.sty}{% use microtype if available
  \usepackage[]{microtype}
  \UseMicrotypeSet[protrusion]{basicmath} % disable protrusion for tt fonts
}{}
\makeatletter
\@ifundefined{KOMAClassName}{% if non-KOMA class
  \IfFileExists{parskip.sty}{%
    \usepackage{parskip}
  }{% else
    \setlength{\parindent}{0pt}
    \setlength{\parskip}{6pt plus 2pt minus 1pt}}
}{% if KOMA class
  \KOMAoptions{parskip=half}}
\makeatother
\usepackage{color}
\usepackage{fancyvrb}
\newcommand{\VerbBar}{|}
\newcommand{\VERB}{\Verb[commandchars=\\\{\}]}
\DefineVerbatimEnvironment{Highlighting}{Verbatim}{commandchars=\\\{\}}
% Add ',fontsize=\small' for more characters per line
\usepackage{framed}
\definecolor{shadecolor}{RGB}{248,248,248}
\newenvironment{Shaded}{\begin{snugshade}}{\end{snugshade}}
\newcommand{\AlertTok}[1]{\textcolor[rgb]{0.94,0.16,0.16}{#1}}
\newcommand{\AnnotationTok}[1]{\textcolor[rgb]{0.56,0.35,0.01}{\textbf{\textit{#1}}}}
\newcommand{\AttributeTok}[1]{\textcolor[rgb]{0.13,0.29,0.53}{#1}}
\newcommand{\BaseNTok}[1]{\textcolor[rgb]{0.00,0.00,0.81}{#1}}
\newcommand{\BuiltInTok}[1]{#1}
\newcommand{\CharTok}[1]{\textcolor[rgb]{0.31,0.60,0.02}{#1}}
\newcommand{\CommentTok}[1]{\textcolor[rgb]{0.56,0.35,0.01}{\textit{#1}}}
\newcommand{\CommentVarTok}[1]{\textcolor[rgb]{0.56,0.35,0.01}{\textbf{\textit{#1}}}}
\newcommand{\ConstantTok}[1]{\textcolor[rgb]{0.56,0.35,0.01}{#1}}
\newcommand{\ControlFlowTok}[1]{\textcolor[rgb]{0.13,0.29,0.53}{\textbf{#1}}}
\newcommand{\DataTypeTok}[1]{\textcolor[rgb]{0.13,0.29,0.53}{#1}}
\newcommand{\DecValTok}[1]{\textcolor[rgb]{0.00,0.00,0.81}{#1}}
\newcommand{\DocumentationTok}[1]{\textcolor[rgb]{0.56,0.35,0.01}{\textbf{\textit{#1}}}}
\newcommand{\ErrorTok}[1]{\textcolor[rgb]{0.64,0.00,0.00}{\textbf{#1}}}
\newcommand{\ExtensionTok}[1]{#1}
\newcommand{\FloatTok}[1]{\textcolor[rgb]{0.00,0.00,0.81}{#1}}
\newcommand{\FunctionTok}[1]{\textcolor[rgb]{0.13,0.29,0.53}{\textbf{#1}}}
\newcommand{\ImportTok}[1]{#1}
\newcommand{\InformationTok}[1]{\textcolor[rgb]{0.56,0.35,0.01}{\textbf{\textit{#1}}}}
\newcommand{\KeywordTok}[1]{\textcolor[rgb]{0.13,0.29,0.53}{\textbf{#1}}}
\newcommand{\NormalTok}[1]{#1}
\newcommand{\OperatorTok}[1]{\textcolor[rgb]{0.81,0.36,0.00}{\textbf{#1}}}
\newcommand{\OtherTok}[1]{\textcolor[rgb]{0.56,0.35,0.01}{#1}}
\newcommand{\PreprocessorTok}[1]{\textcolor[rgb]{0.56,0.35,0.01}{\textit{#1}}}
\newcommand{\RegionMarkerTok}[1]{#1}
\newcommand{\SpecialCharTok}[1]{\textcolor[rgb]{0.81,0.36,0.00}{\textbf{#1}}}
\newcommand{\SpecialStringTok}[1]{\textcolor[rgb]{0.31,0.60,0.02}{#1}}
\newcommand{\StringTok}[1]{\textcolor[rgb]{0.31,0.60,0.02}{#1}}
\newcommand{\VariableTok}[1]{\textcolor[rgb]{0.00,0.00,0.00}{#1}}
\newcommand{\VerbatimStringTok}[1]{\textcolor[rgb]{0.31,0.60,0.02}{#1}}
\newcommand{\WarningTok}[1]{\textcolor[rgb]{0.56,0.35,0.01}{\textbf{\textit{#1}}}}
\usepackage{graphicx}
\makeatletter
\newsavebox\pandoc@box
\newcommand*\pandocbounded[1]{% scales image to fit in text height/width
  \sbox\pandoc@box{#1}%
  \Gscale@div\@tempa{\textheight}{\dimexpr\ht\pandoc@box+\dp\pandoc@box\relax}%
  \Gscale@div\@tempb{\linewidth}{\wd\pandoc@box}%
  \ifdim\@tempb\p@<\@tempa\p@\let\@tempa\@tempb\fi% select the smaller of both
  \ifdim\@tempa\p@<\p@\scalebox{\@tempa}{\usebox\pandoc@box}%
  \else\usebox{\pandoc@box}%
  \fi%
}
% Set default figure placement to htbp
\def\fps@figure{htbp}
\makeatother
\setlength{\emergencystretch}{3em} % prevent overfull lines
\providecommand{\tightlist}{%
  \setlength{\itemsep}{0pt}\setlength{\parskip}{0pt}}
\usepackage{bookmark}
\IfFileExists{xurl.sty}{\usepackage{xurl}}{} % add URL line breaks if available
\urlstyle{same}
\hypersetup{
  pdftitle={Survival primer: KM{,} log-rank{,} Cox PH},
  hidelinks,
  pdfcreator={LaTeX via pandoc}}

\title{Survival primer: KM, log-rank, Cox PH}
\author{}
\date{\vspace{-2.5em}}

\begin{document}
\maketitle

\textbf{Inference lab:} Hypothesis → Visualise → Assumptions → Conduct →
Conclude.

\subsection{Learning objectives}\label{learning-objectives}

\begin{itemize}
\tightlist
\item
  Estimate and plot a Kaplan-Meier survival curve for a two-group
  comparison.
\item
  Run and interpret a log-rank test.
\item
  Fit a Cox proportional-hazards model, check the PH assumption, and
  read a hazard ratio honestly.
\end{itemize}

\subsection{Prerequisites}\label{prerequisites}

Comfort with GLMs and with reading a regression table.

\subsection{Background}\label{background}

Survival analysis handles outcomes that are times to an event, where
some subjects have not yet had the event when observation ends --- the
familiar right-censoring problem. Ignoring censoring by throwing out the
censored observations biases the estimator; pretending the censored
observations are event-free biases it the other way. The survival
machinery handles censoring correctly by reasoning at each event time
about who is still at risk.

Three tools cover most clinical applications: the Kaplan-Meier estimator
for a non-parametric survival curve, the log-rank test for a two-group
comparison, and Cox proportional-hazards regression for adjustment and
multivariable modelling. The PH assumption --- that the hazard ratio is
constant over time --- deserves explicit checking; when it fails, a
time-varying coefficient or a stratified model is usually the right fix.

\subsection{Setup}\label{setup}

\begin{Shaded}
\begin{Highlighting}[]
\FunctionTok{library}\NormalTok{(tidyverse)}
\FunctionTok{library}\NormalTok{(survival)}
\FunctionTok{library}\NormalTok{(ggsurvfit)}
\FunctionTok{library}\NormalTok{(broom)}
\FunctionTok{set.seed}\NormalTok{(}\DecValTok{42}\NormalTok{)}
\FunctionTok{theme\_set}\NormalTok{(}\FunctionTok{theme\_minimal}\NormalTok{(}\AttributeTok{base\_size =} \DecValTok{12}\NormalTok{))}
\end{Highlighting}
\end{Shaded}

\subsection{1. Hypothesis}\label{hypothesis}

Using \texttt{survival::lung}, compare survival between male and female
patients. H₀: hazard ratio = 1. H₁: hazard ratio ≠ 1. α = 0.05.

\subsection{2. Visualise}\label{visualise}

\begin{Shaded}
\begin{Highlighting}[]
\NormalTok{lung2 }\OtherTok{\textless{}{-}}\NormalTok{ lung }\SpecialCharTok{|\textgreater{}}
  \FunctionTok{mutate}\NormalTok{(}\AttributeTok{sex =} \FunctionTok{factor}\NormalTok{(sex, }\AttributeTok{levels =} \DecValTok{1}\SpecialCharTok{:}\DecValTok{2}\NormalTok{, }\AttributeTok{labels =} \FunctionTok{c}\NormalTok{(}\StringTok{"male"}\NormalTok{, }\StringTok{"female"}\NormalTok{)))}

\NormalTok{km }\OtherTok{\textless{}{-}} \FunctionTok{survfit2}\NormalTok{(}\FunctionTok{Surv}\NormalTok{(time, status) }\SpecialCharTok{\textasciitilde{}}\NormalTok{ sex, }\AttributeTok{data =}\NormalTok{ lung2)}

\NormalTok{km }\SpecialCharTok{|\textgreater{}}
  \FunctionTok{ggsurvfit}\NormalTok{() }\SpecialCharTok{+}
  \FunctionTok{add\_confidence\_interval}\NormalTok{() }\SpecialCharTok{+}
  \FunctionTok{add\_risktable}\NormalTok{() }\SpecialCharTok{+}
  \FunctionTok{labs}\NormalTok{(}\AttributeTok{x =} \StringTok{"Days"}\NormalTok{, }\AttributeTok{y =} \StringTok{"Survival probability"}\NormalTok{)}
\end{Highlighting}
\end{Shaded}

\subsection{3. Assumptions}\label{assumptions}

Two assumptions are worth checking: non-informative censoring (patients
censored at a given time are representative of those still at risk) and
proportional hazards.

\begin{Shaded}
\begin{Highlighting}[]
\FunctionTok{cox.zph}\NormalTok{(cox\_fit)}
\end{Highlighting}
\end{Shaded}

A large \emph{p}-value on each row and a non-systematic Schoenfeld
residual plot support the PH assumption; a small \emph{p} calls for a
time-varying coefficient or stratification.

\subsection{4. Conduct}\label{conduct}

\begin{Shaded}
\begin{Highlighting}[]
\NormalTok{lr}

\FunctionTok{tidy}\NormalTok{(cox\_fit, }\AttributeTok{exponentiate =} \ConstantTok{TRUE}\NormalTok{, }\AttributeTok{conf.int =} \ConstantTok{TRUE}\NormalTok{)}
\end{Highlighting}
\end{Shaded}

\subsection{5. Concluding statement}\label{concluding-statement}

\begin{Shaded}
\begin{Highlighting}[]
\NormalTok{hr\_row }\OtherTok{\textless{}{-}} \FunctionTok{tidy}\NormalTok{(cox\_fit, }\AttributeTok{exponentiate =} \ConstantTok{TRUE}\NormalTok{, }\AttributeTok{conf.int =} \ConstantTok{TRUE}\NormalTok{) }\SpecialCharTok{|\textgreater{}}
  \FunctionTok{filter}\NormalTok{(term }\SpecialCharTok{==} \StringTok{"sexfemale"}\NormalTok{)}
\end{Highlighting}
\end{Shaded}

\begin{quote}
Female sex was associated with a lower risk of death relative to male
sex (HR \texttt{round(hr\_row\$estimate,\ 2)}; 95\% CI
\texttt{round(hr\_row\$conf.low,\ 2)} to
\texttt{round(hr\_row\$conf.high,\ 2)}; \emph{p} =
\texttt{signif(hr\_row\$p.value,\ 2)}), adjusted for age. The
proportional-hazards assumption was not clearly violated (Schoenfeld
test, all \emph{p} \textgreater{} 0.05).
\end{quote}

\subsection{Common pitfalls}\label{common-pitfalls}

\begin{itemize}
\tightlist
\item
  Reporting a median survival that is not reached and calling it
  ``undefined'' in the results. Say so explicitly.
\item
  Ignoring the PH check because the output looks clean.
\item
  Comparing Kaplan-Meier curves by eye alone --- add the log-rank test
  and the risk table.
\item
  Collapsing a competing event (e.g.~death from another cause) into
  censoring without justification.
\end{itemize}

\subsection{Further reading}\label{further-reading}

\begin{itemize}
\tightlist
\item
  Harrell FE. \emph{Regression Modeling Strategies}, ch.~20.
\item
  Therneau TM \& Grambsch PM. \emph{Modeling Survival Data.}
\end{itemize}

\subsection{Session info}\label{session-info}

\begin{Shaded}
\begin{Highlighting}[]

\end{Highlighting}
\end{Shaded}

\subsection{See also --- chapter index}\label{see-also-chapter-index}

\begin{itemize}
\tightlist
\item
  \href{../index.Rmd}{Methods in this chapter}
\item
  \href{../../../ch00-find-your-method.Rmd}{Find your method (Chapter
  0)}
\end{itemize}

\end{document}
